% ============================================
% DOC-05 (ZH)
% 中文翻译 — SPEC-935-R356
% ============================================
\section*{\textbf{DOC-05} --- 会谈报告II}
\markboth{DOC-05 --- 会谈报告II}{}

\noindent \textit{埃托尔·奥利维拉医生 —— 1979年3月21日}

\noindent \textbf{埃托尔·奥利维拉医生 —— 1979年3月21日}

\vspace{0.5em}\noindent\rule{\textwidth}{0.5pt}\vspace{0.5em}

\noindent \textbf{患者:} 若阿金·安东尼奥·科雷亚 —— 患者1,260
\textbf{类型:} 精神科评估会谈
\textbf{时长:} 90分钟（延长）

\vspace{0.5em}\noindent\rule{\textwidth}{0.5pt}\vspace{0.5em}

\noindent \textbf{患者观察:}

\noindent 患者今天表现不同。更兴奋。近乎欣快。我进房间时他在微笑——一个不对称的微笑，向脸部右侧牵拉，那里黄眼比上次会谈时更亮。

\noindent 他问：\textit{“您做梦了吗？”}

\noindent 我回答做了，做了一个不寻常的梦。他没有要求细节。只是点了点头，仿佛他已经知道。

\noindent 本次会谈中，若阿金连续说了75分钟。几乎不间断。他描述了在塞阿拉的童年、1915年的大旱、从塞纳多尔·蓬佩乌向福塔莱萨的跋涉、父亲的死，以及在阿拉加迪索（Alagadiço）收容营的关押经历。叙述生动——有感官细节（气味、温度、质感），即使在记忆保存完好的患者身上也不常见。

\noindent 当他描述“相遇”时，叙述的语调变了。

\noindent 据若阿金说，他在集中营后面的尸体坑里遇到了一只怪物。描述包含精神病性要素（黄眼、骨刺、外露的肋骨），但\textit{内在一致性}令人瞩目。他描述这次相遇时的确信程度，与描述父亲之死时相同。

\noindent \textbf{逐字记录:}

\noindent \textit{“它说：‘你也很饿。’我说是的。它说：‘我替你吃。我替所有吃不下的人吃。’于是我坐下来。坐在坑边，看着它吃。我不感到害怕。我感到饥饿。一种再也没有消失的饥饿。”}

\noindent \textbf{分析:}

\noindent - “怪物”似乎是若阿金童年创伤的表现，可能是他在童年目睹的饥饿与死亡的化身。
- 坑边的相遇可能代表一个解离时刻：若阿金可能目睹了无法承受的东西（尸体坑、他人的死亡），他的头脑创造了一个替代性叙事来处理该事件。
- “永不消失的饥饿”与创伤后应激障碍的临床相相符，在创伤后应激障碍中，创伤事件被无限地反复体验。

\noindent \textbf{治疗计划:}

\noindent - 我申请阅读患者的日记，起初被拒绝，但患者同意“等准备好时”交给我。
- 维持每周会谈。

\noindent \textbf{个人注:}

\noindent 我开始做反复出现的梦。不是普通的噩梦——极其生动，带有气味和触觉感受。我醒来时嘴里有泥土的味道。我没有对若阿金提起，但他知道。他知道。

\vspace{0.5em}\noindent\rule{\textwidth}{0.5pt}\vspace{0.5em}

\noindent \textit{埃托尔·奥利维拉医生}
\textit{精神科医生 —— 最高安全病区}
\textit{巴尔巴塞纳收容医院}

\vspace{0.5em}
\noindent\textit{\textbf{↪ 路线:} \rota{MEM-15} \rotasep \rota{DOC-07} \rotasep \rota{LUC-03}}
